\documentclass[11pt]{article}
\usepackage[margin=1in]{geometry}
\usepackage{amsmath,amssymb,amsthm}
\usepackage{mathtools}
\usepackage[numbers,square]{natbib}
\usepackage{booktabs}
\usepackage{csquotes}
\usepackage{hyperref}
\hypersetup{hypertexnames=false}
\usepackage{cleveref}

% Theorem environments (matching the ecosystem preamble)
\theoremstyle{definition}
\newtheorem{definition}{Definition}[section]
\theoremstyle{plain}
\newtheorem{theorem}{Theorem}[section]
\newtheorem{proposition}[theorem]{Proposition}
\newtheorem{corollary}[theorem]{Corollary}
\newtheorem{lemma}[theorem]{Lemma}
\newtheorem{conjecture}[theorem]{Conjecture}
\theoremstyle{definition}
\newtheorem{assumption}{Assumption}[section]
\theoremstyle{remark}
\newtheorem{remark}{Remark}[section]
\newtheorem{example}{Example}[section]

% Shared macros (ecosystem-canonical; see formalism/notation.md)
\newcommand{\enc}{\mathrm{enc}}
\newcommand{\dec}{\mathrm{dec}}
\newcommand{\fhat}{\hat{f}}
\newcommand{\ghat}{\hat{g}}
\newcommand{\cipher}[1]{\mathsf{C}(#1)}
\newcommand{\cipherS}[2]{\mathsf{C}_{#2}(#1)}
\newcommand{\orbitF}{\mathrm{orbit}_F}
\newcommand{\rekey}[2]{r_{#1 \to #2}}
\newcommand{\B}{\{0,1\}}

\title{Cipher Closures: Cipher Data Structures and the\\
Code--Data Duality in Trapdoor Computing}
\author{Alexander Towell\\\texttt{lex@metafunctor.com}}
\date{\today}

\begin{document}
\maketitle

\begin{abstract}
A cipher value is opaque data; a cipher map is opaque code. We show
that in trapdoor computing the two are interchangeable, and that the
cipher map is a universal carrier for both. A cipher data structure
(a pair, a list, a map) is realized as a single cipher map: its
operation interface is folded into a query tag (the \emph{dispatch
pattern}), and sequential structure is encoded in the shape of an
orbit under a hash chain (\emph{designed orbits}). Conversely, a
cipher map is itself a point of a cipher type, the cipher exponential
$\cipher{A \to B}$, so code is data. The abstraction that makes the
duality precise is the \emph{cipher closure}: a procedure that
captures a secret and exposes an operation interface, of which cipher
values, cipher maps, and cipher data structures are special cases.
Cipher rekeying, which transforms the captured secret, is then one
instance of a general principle: in trapdoor computing, values,
functions, data structures, and the secret itself are all cipher
data, carried by cipher maps. We give the constructions, their
information-theoretic leakage (orbit-encoded structures leak length
and iteration order but bound per-element confidentiality by the
orbit-closure bound), and the unification.
\end{abstract}

%% ====================================================================
\section{Introduction}
\label{sec:intro}
%% ====================================================================

In the Lisp tradition, code and data are made of the same stuff. A
program is an S-expression, an S-expression is a list, and a list is a
value the language can build, inspect, and run~\cite{mccarthy1960recursive}.
The clearest small instance is the closure. In Abelson and Sussman's
exposition~\cite{abelson1996sicp}, a pair needs no primitive storage:
it is a procedure that has captured two values and dispatches on the
operation asked of it,
\[
\mathrm{cons}(x, y) = \lambda m.\, m(x, y),\qquad
\mathrm{car}(p) = p(\lambda x\, y.\, x),\qquad
\mathrm{cdr}(p) = p(\lambda x\, y.\, y).
\]
The pair is data (it holds $x$ and $y$) and code (it is a procedure)
at once. This is the move that lets one speak of a data structure as a
function and a function as a value, the homoiconicity of the
$\lambda$-calculus and its descendants~\cite{landin1964mechanical}.

Trapdoor computing~\cite{towell2026cipher} supplies a single carrier
that makes the same duality concrete over a secret. A \emph{cipher
value} $c \in \B^n$ is an opaque bit string that encodes a latent
value $x$ under a secret $s$, decodable only by a holder of $s$. A
\emph{cipher map} $\fhat : \B^n \to \B^n$ is an opaque total function
that realizes a latent function $f : X \to Y$: an untrusted machine
can evaluate it without learning $x$, $f$, or $s$. The untrusted
machine sees bit strings flowing through a lookup table; the trusted
machine holds the encoder, the decoder, and the secret. Cipher values
are opaque data; cipher maps are opaque code. The question of this
paper is how far the boundary between the two dissolves, and what is
gained when it does.

We answer it in three directions, organized around one object, the
\emph{cipher closure}: a procedure that captures a secret and exposes
an operation interface, of which cipher values, cipher maps, and
cipher data structures are special cases.

\paragraph{Data is code.} A cipher data structure with several
operations need not be a heap of storage with a protocol bolted on. It
is a single cipher map. The operation interface folds into a query
tag, so one tagged cipher map answers \texttt{is\_member} and
\texttt{lookup} alike (\Cref{sec:dispatch}, the \emph{dispatch
pattern}); and sequential structure folds into the shape of an orbit
under a hash chain, so one cipher map and a starting key encode a list,
a map, or a vector (\Cref{sec:designed-orbits}, \emph{designed
orbits}). Algebraic cipher types~\cite{towell2026algebraic} reads an
orbit as a leakage measure: the more cipher values an adversary can
reach by composing the operations it holds, the more it can learn. We
turn that reading around. Shape the orbit deliberately and it is no
longer only a leakage budget to be bounded; it is a construction
primitive, and the same bound that limited the adversary now measures
exactly what an orbit-encoded data structure reveals
(\Cref{sec:leakage}).

\paragraph{Code is data.} The reverse direction is already a theorem
of the type algebra. The cipher exponential $\cipher{A \to B}$ is
\emph{precisely} the set of cipher maps $\cipher{A} \to
\cipher{B}$~\cite[Prop.~4.3]{towell2026algebraic}. So a cipher map
$\fhat$ is at once code (the realization of a latent function) and a
point of a cipher type (an inhabitant of $\cipher{A \to B}$). Combined
with the first direction, the cipher map carries both: a data
structure becomes a cipher map, and a cipher map is a cipher value of
function type (\Cref{sec:exponential}).

\paragraph{The secret is data.} The captured secret is itself state in
a closure, and transforming it is one more instance of the same move.
Cipher rekeying~\cite{towell2026rekeying} re-parameterizes the secret
a cipher object captured, taking a cipher value under one secret to a
cipher value of the same latent value under another. It is the
secret-as-data corner of the duality, and, with the cipher exponential,
one of its two already-charted endpoints (\Cref{sec:duality}).

\paragraph{Contributions.}
\begin{enumerate}
\item \textbf{The code--data duality for trapdoor computing}, made
  precise through cipher closures (\Cref{sec:closures}): homoiconicity
  over a secret, with cipher values, cipher maps, and cipher data
  structures as the three subclasses of one object.
\item \textbf{The dispatch pattern} (\Cref{sec:dispatch}): a
  multi-operation cipher data structure realized as a single tagged
  cipher map, the same boundary that
  \cite{towell2026cipherprog,towell2026algebraic} cut through an
  expression tree, drawn here from the data-structure side.
\item \textbf{Designed orbits} (\Cref{sec:designed-orbits}): XOR-orbit
  hash chains that encode lists, two-level maps, and indexed vectors
  inside a single cipher map, inverting the orbit-as-leakage reading
  of~\cite{towell2026algebraic} into an orbit-as-construction
  primitive, with one cipher map yielding many data structures as the
  starting key varies.
\item \textbf{An information-theoretic leakage analysis} of
  orbit-encoded data structures (\Cref{sec:leakage}): the list length
  and the iteration order leak, while per-element confidentiality is
  bounded by the orbit-closure bound; what is provable is proved, and
  what is not is marked as a conjecture or open problem.
\item \textbf{The unification} (\Cref{sec:duality}): values, functions,
  data structures, and the secret are all cipher data carried by cipher
  maps, with the cipher exponential and cipher rekeying as the two
  known endpoints.
\end{enumerate}

\paragraph{Position.} The foundation
paper~\cite{towell2026cipher} establishes the cipher map and treats it
as code (a function $X \to Y$); it does not consider data structures.
Algebraic cipher types~\cite{towell2026algebraic} supplies the
exponential identity and the orbit-closure bound, but reads orbits only
as a leakage measure. Cipher rekeying~\cite{towell2026rekeying}
develops the secret-as-data instance and explicitly defers the
general duality to this paper. A sibling
paper~\cite{towell2026cipherprog} realizes whole \emph{programs} as
cipher-map compositions by cutting an expression tree; this paper is
about cipher \emph{data structures} and the code--data frame, an
orthogonal concern that meets program construction only at the shared
cut boundary of \Cref{sec:dispatch}.

\paragraph{Scope.} The privacy model throughout is
information-theoretic, resting on a one-way hash, uniform
representation, and the finiteness of an orbit; we import no
ORAM, simulation-based, or game-based cryptographic machinery. The
categorical language (functor, exponential object) is an organizing
device that makes the notation regular; it is not load-bearing in any
proof, and we flag where it is doing only bookkeeping.

%% ====================================================================
\section{Preliminaries}
\label{sec:prelim}
%% ====================================================================

We recall the cipher map abstraction and the two results from the
type algebra that this paper builds on, then fix the closure notation.
Throughout, cipher values are $n$-bit strings over $\B = \{0,1\}$. We
use entropy, conditional entropy, mutual information, and total
variation distance in their standard senses without redefining them.

\subsection{Cipher maps and cipher spaces}

\begin{definition}[Cipher map~{\cite[Def.~3.1]{towell2026cipher}}]
\label{def:cipher-map}
Let $f : X \to Y$ be a \emph{latent function}. A \emph{cipher map} for
$f$ is a tuple $(\fhat, \enc, \dec, s)$ where
$\fhat : \B^n \to \B^n$ is a total function on $n$-bit strings;
$\enc : X \times \{0, \ldots, K(x){-}1\} \to \B^n$ is an encoding
function, with $K(x) \geq 1$ the number of encodings of $x$;
$\dec : \B^n \to Y \cup \{\bot\}$ is a decoding function; and $s$ is
the secret (the trapdoor) from which $\fhat$, $\enc$, and $\dec$ are
derived. The untrusted machine holds $\fhat$; the trusted machine
holds $\enc$, $\dec$, and $s$.
\end{definition}

A cipher map satisfies four properties parameterized by
$(\eta, \varepsilon, \mu, \delta)$~\cite[Sec.~4]{towell2026cipher},
which we use but do not re-derive:
\emph{totality} ($\fhat$ is defined on all of $\B^n$, and on inputs
outside the encoded domain its output is uniform under the
random-oracle model, so the untrusted machine cannot tell a real query
from filler); \emph{representation uniformity} (the distribution of
cipher values the untrusted machine observes is within total variation
$\delta$ of uniform); \emph{correctness} ($\dec(\fhat(\enc(x,k))) =
f(x)$ except with probability at most $\eta$); and \emph{composability}
(the composition $\ghat \circ \fhat$ is a cipher map for $g \circ f$,
with errors combining as $\eta_{g\circ f} = 1 - (1-\eta_f)(1-\eta_g)$).

We write $\cipherS{\cdot}{s}$ for the \emph{cipher functor} at secret
$s$, dropping the subscript to $\cipher{\cdot}$ when $s$ is fixed. On a
latent type $X$ it returns the \emph{cipher space} $\cipherS{X}{s}
\subseteq \B^n$, the set of bit strings encoding elements of $X$ under
$s$; on a latent function $f$ it returns the cipher map $\fhat$, so
$\cipherS{f}{s} = \fhat$. A cipher value of $x$ is $\enc_s(x, k)$ for
some representation index $k$; the hat marks a cipher \emph{map}
(the only approximate object), never a value or a secret.

\subsection{The orbit-closure bound}

The single tool we borrow from the type algebra is a bound on what an
active untrusted machine can learn by composing the cipher maps it
holds.

\begin{definition}[Orbit closure~{\cite[Def.~5.1]{towell2026algebraic}}]
\label{def:orbit}
Let $F = \{\fhat_1, \ldots, \fhat_k\}$ be cipher maps available to the
untrusted machine and $c \in \B^n$ a cipher value. The \emph{orbit
closure} $\orbitF(c) \subseteq \B^n$ is the smallest set containing
$c$ and closed under every $\fhat_i \in F$: the set of all bit strings
reachable from $c$ by composing operations in $F$ in any order. Since
$\B^n$ is finite, $\orbitF(c)$ is finite, and it depends only on $F$
and $c$, not on the latent function or the trapdoor.
\end{definition}

\begin{theorem}[Orbit-closure bound~{\cite[Thm.~5.3]{towell2026algebraic}}]
\label{thm:orbit-bound}
Let $X$ be the latent value encoded by $c$, drawn from the trusted
machine's input distribution, and let the untrusted machine's
observation after probing $c$ with $F$ be a random variable supported
on $\orbitF(c)$. Then
\[
H(X \mid \mathcal{V}_F(c)) \;\geq\; H(X) - \log_2 |\orbitF(c)|.
\]
\end{theorem}

The proof is a one-line counting bound: the observation takes at most
$|\orbitF(c)|$ values, so it carries at most $\log_2 |\orbitF(c)|$ bits
about $X$. The content is in identifying the orbit as the support of
what the adversary sees. The bound is an \emph{active}-adversary
statement on a different axis from the marginal parameter $\delta$:
$\delta$ governs the leakage of a single cipher value at rest, the
orbit governs the leakage of an operation interface under repeated
use, and neither bounds the other. It is loose (it ignores the
functional relations among orbit elements) and is approached with
equality when the cipher maps behave as independent random
oracles~\cite{bellare1993random}. These are exactly the properties we
will need in \Cref{sec:leakage}.

\subsection{Cipher closures}

A closure, in the SICP sense~\cite{abelson1996sicp}, is a procedure
that has captured values from its environment. We work throughout with
closures that capture a secret and expose an operation interface.

\begin{definition}[Cipher closure]
\label{def:cipher-closure}
A \emph{cipher closure} over secret $s$ is a function $\kappa_s :
\mathcal{O} \to \mathcal{R}$ where $\mathcal{O}$ is a set of operations
and $\mathcal{R}$ the union of their result types. The closure has
access to $s$ and to any auxiliary state captured at construction time;
each operation $o \in \mathcal{O}$ has its own input and output types,
which the abstract signature suppresses.
\end{definition}

%% ====================================================================
\section{Cipher Closures}
\label{sec:closures}
%% ====================================================================

The cipher closure is the homoiconic object of trapdoor computing: it
is data (the state it captured, including the secret) and code (the
procedure it is) at the same time. The three primitives of the
framework are its three subclasses, distinguished by what they expose
and, crucially, by who can run them.

\begin{description}
\item[Cipher value closures.] A cipher value $\enc_s(x, k)$ is the
  closure that captured the bit string and the secret and exposes a
  single operation, $\texttt{decode}$, which returns the latent value.
  Decoding requires the secret, so this operation is invocable only by
  the trusted machine; to the untrusted machine a cipher value is inert
  data, a bit string with no working interface. In the cipher functor
  view it is a point of a cipher space, $\enc_s(x,k) \in \cipherS{X}{s}$.

\item[Cipher map closures.] A cipher map $\fhat$ is the closure that
  captured the secret and exposes a single operation,
  \emph{evaluation}: apply $\fhat$ to a cipher value to get a cipher
  value. This is the one subclass the untrusted machine can run, and
  it is exactly what the untrusted machine holds
  (\Cref{def:cipher-map}). Evaluation reveals nothing about the secret
  because totality makes every input admissible and every output
  $n$ bits wide. A cipher map closure is at once code (it realizes a
  latent function) and, by the cipher exponential
  identity~\cite[Prop.~4.3]{towell2026algebraic}, a point of the cipher
  type $\cipher{A \to B}$; \Cref{sec:exponential} returns to this.

\item[Cipher data structure closures.] A cipher data structure
  captures internal state, possibly mutable, and exposes several
  operations, each realized by a cipher map (so each is runnable by the
  untrusted machine) or by a captured-value return (so it is
  operationally free but visible). This is the subclass the rest of the
  paper develops: \Cref{sec:dispatch} collapses its several operations
  into one tagged cipher map, and \Cref{sec:designed-orbits} encodes
  its sequential structure in an orbit.
\end{description}

The boundary between the subclasses is the boundary between data and
code, and the closure straddles it. The next three examples are the
classical data structures realized as cipher closures. We state each
construction self-contained, and after each we state honestly what it
leaks; the leakage caveats are not afterthoughts but the substance of
\Cref{sec:leakage}.

\begin{example}[Cipher Boolean]
\label{ex:cipher-bool}
The cipher Boolean type partitions $\B^n$ into True, False, and Noise
regions fixed by the secret~\cite[Sec.~7.1]{towell2026algebraic}. A
cipher Boolean value is the closure that captures the secret and a bit
string and exposes $\texttt{decode}() \to \mathrm{Bool}$ (trusted-only,
it reads the partition), $\texttt{and}(b') \to \cipher{\mathrm{Bool}}$,
and $\texttt{not}() \to \cipher{\mathrm{Bool}}$. The latter two are
realized by cipher maps and are exposed to the untrusted machine. They
are not free of consequence: an adversary that computes
$\texttt{and}(b, \texttt{not}(b))$ obtains an (approximate) cipher
encoding of False, which it can use to sort other cipher Booleans by
value~\cite[Ex.~5.1]{towell2026algebraic}. This is orbit-closure
leakage of the operation interface, bounded by \Cref{thm:orbit-bound}.
\end{example}

\begin{example}[Cipher pair via cons/car/cdr]
\label{ex:cipher-pair}
Following Abelson and Sussman, a cipher pair of cipher values $a \in
\cipher{A}$ and $b \in \cipher{B}$ is the closure
\[
\mathrm{cipher\_cons}(a, b) = \lambda m.\, m(a, b),\qquad
\mathrm{cipher\_first}(p) = p(\lambda a\, b.\, a),\qquad
\mathrm{cipher\_second}(p) = p(\lambda a\, b.\, b).
\]
This realizes the \emph{componentwise} product $\cipher{A} \times
\cipher{B}$ of~\cite[Sec.~4.2]{towell2026algebraic}: projection is
operationally free, since the closure returns its captured component
and no cipher map is invoked. It is not leakage-free. The adversary
now holds both cipher values from one closure, and the empirical joint
distribution of cipher pairs across many closures converges to the
joint distribution of the latent pairs
$(a, b)$~\cite[Prop.~4.1]{towell2026algebraic}. The alternative, the
\emph{joint} product $\cipher{A \times B}$, encodes the pair as a
single cipher value and hides the correlation, at the cost of
cipher-map-mediated projections rather than free ones. This is the
encoding-granularity trade-off of~\cite[Sec.~8.1]{towell2026cipher},
made visible in closure form: componentwise captures two cipher
values, joint captures one plus the projection maps.
\end{example}

\begin{example}[Mutable cipher counter]
\label{ex:cipher-counter}
A closure can capture mutable state. A cipher counter at $0$ is
\[
\mathrm{cipher\_counter}() = \texttt{let}\; c = \enc(0, 0)\;
\texttt{in}\; \lambda \texttt{op}.\,
\begin{cases}
c & \texttt{op} = \texttt{read},\\
c \leftarrow \widehat{\mathrm{succ}}(c);\, c & \texttt{op} = \texttt{increment},
\end{cases}
\]
with $\widehat{\mathrm{succ}}$ the cipher successor map. The closure
mutates the captured cipher value in place, so closures are not
restricted to pure functions: with captured mutable state they realize
mutable cipher data structures. The counter leaks its orbit. An
adversary with active access to the operations and a known starting
state (zero is plaintext) calls $\texttt{increment}$ from the known
$c_0$ and recovers $\{c_0, \widehat{\mathrm{succ}}(c_0),
\widehat{\mathrm{succ}}^2(c_0), \ldots\}$, pairing each cipher value
with its plaintext index. This is precisely the successor leakage
of~\cite[Ex.~5.2]{towell2026algebraic}, bounded by
\Cref{thm:orbit-bound}; the closure neither weakens nor strengthens
it, only packages the operations through which the orbit is reached.
Standard mitigations apply (randomize the start, use multiple
representations so $\widehat{\mathrm{succ}}$ fans out, or cap the orbit
with typed composition chains); none affects mutability.
\end{example}

These examples make the same point at three difficulties: a cipher
value, a cipher map, and a cipher data structure are all cipher
closures over a secret, and each is data and code at once. What
remains is to show how a data structure becomes a single cipher map
(the two directions of \enquote{data is code},
\Cref{sec:dispatch,sec:designed-orbits}), to read off what those
constructions leak (\Cref{sec:leakage}), and to recall the reverse
direction (\enquote{code is data}, \Cref{sec:exponential}).

%% ====================================================================
\section{Data is Code: The Dispatch Pattern}
\label{sec:dispatch}
%% ====================================================================

A cipher data structure with several operations is realizable as a
single cipher map whose first input is an operation tag. Take a cipher
list, with a query type $\mathcal{Q}$ tagging the requested operation:
\[
\fhat : \mathcal{Q} \times \cipher{X} \to \cipher{Y \cup \mathrm{Bool}},
\qquad
\fhat(\texttt{key\_query}, k) = \begin{cases}
\cipher{\mathrm{True}} & k \in \mathrm{keys}(\fhat)\\
\cipher{\mathrm{False}} & \text{otherwise}
\end{cases},
\]
\[
\fhat(\texttt{map\_query}, k) = \cipher{\mathrm{value}\;\mathrm{at}\;k}.
\]
One cipher map realizes both \texttt{is\_member} and \texttt{lookup}
semantics; the trusted machine constructs one cipher map per closure
rather than one per operation, and the operations dispatch on the tag
at evaluation time.

The construction generalizes. A cipher data structure closure exposing
operations $o_1, \ldots, o_m$, each $o_j$ realized by a cipher map
$\fhat_j : \cipher{A_j} \to \cipher{B_j}$, is realized by the single
cipher map
\[
\fhat : \mathcal{Q} \times \cipher{A} \to \cipher{B},
\qquad
\fhat(\langle j \rangle, c) = \fhat_j(c),
\]
where $\mathcal{Q}$ is a tag type with one tag $\langle j \rangle$ per
operation, $\cipher{A} = \bigcup_j \cipher{A_j}$ is the union of the
operand spaces, and $\cipher{B} = \bigcup_j \cipher{B_j}$ the union of
the result spaces. The tag is the first coordinate of the input; the
trusted machine builds one cipher map for the whole closure rather
than one per operation, and the operations dispatch on the tag at
evaluation time. Two qualifications matter. First, the tag is data the
untrusted machine supplies and therefore sees: dispatch hides
\emph{which latent value} an operation is applied to, not \emph{which
operation} is being asked, exactly as a method name is visible in an
ordinary object. Hiding the choice of operation is a separate problem
(it is a cipher \emph{sum} over the operations, and the cipher Boolean
that selects among branches is itself a value to be encoded), which we
do not solve here. Second, folding many operations into one cipher map
enlarges that map's domain and hence the orbit an active adversary can
explore; the dispatch pattern trades a smaller construction (one
object instead of $m$) against a larger reachable set, and
\Cref{thm:orbit-bound} prices the trade.

\paragraph{The same boundary, drawn from the other side.} The
tag-folding here is the data-structure-side view of a boundary that a
companion line of work approaches from the program side. Realizing a
\emph{program} as cipher maps cuts its expression tree at chosen nodes:
each \enquote{cut point} marks a subtree replaced by a cipher map, with
plaintext above the cut on the trusted machine and cipher values
below~\cite{towell2026cipherprog}. A cipher data structure is the
limiting case where the whole interface is below a single cut: its
operations are the subtrees, and the dispatch tag is the selector that
chooses which subtree to run. The two papers therefore meet at one
construct seen two ways. \cite{towell2026cipherprog} asks, given a
fixed program, where to place cuts so that what stays plaintext leaks
little and what becomes cipher costs little; we ask, given a data
structure, how to present its operations as one cipher map. The
control-flow obliviousness analysis there (branching is untrusted
dispatch on a sum type, with no oblivious short-circuit) is precisely
the first qualification above, transported to programs.

\paragraph{Properties carry over.} Because $\fhat$ is built by the
same trusted-machine construction as any cipher map, the four
properties of \Cref{def:cipher-map} hold of the dispatched map
uniformly. Totality holds because $\fhat$ is total on $\mathcal{Q}
\times \B^n$ (out-of-domain tag/value pairs map to noise like any other
filler). Representation uniformity and correctness are inherited
per-operation from the $\fhat_j$, with $\delta = \max_j \delta_j$ and
$\eta = \max_j \eta_j$ over the operand distributions, since an
adversary query lands in exactly one operation's region. Composability
is unaffected: a dispatched map composes with another cipher map on the
result coordinate exactly as its constituents would. The dispatch
pattern thus costs nothing in the property budget; it only relocates
$m$ separate maps into one, which is what makes \enquote{a data
structure is a cipher map} literally true and not merely a slogan.

%% ====================================================================
\section{Designed Orbits: Data Structures via Hash Chains}
\label{sec:designed-orbits}
%% ====================================================================

Where algebraic-cipher-types reads an orbit as a \emph{leakage}
measure, here we read it as a \emph{constructive} primitive: shape the
orbit and you have encoded a data structure inside a single cipher
map. A list $[y_0, y_1, y_2, \ldots]$ is encoded by one cipher map
$\fhat$ and a starting key $k$ via an XOR hash chain,
\[
y_0 = \fhat(k), \quad
y_1 = \fhat\!\big(k \oplus h(y_0)\big), \quad
y_2 = \fhat\!\big(k \oplus h(y_0) \oplus h(y_1)\big), \quad \ldots
\]
terminating at a sentinel. The iterator advance is
$h \mathrel{\oplus}= h(y); \; y = \fhat(h)$: the accumulator $h$
depends on the entire history of values seen so far, so the untrusted
machine sees a cipher map applied to a sequence of keys that appear
unrelated, with no visible iteration order.

Let us state the list construction precisely. Fix a cryptographic hash
$h : \B^n \to \B^n$ and a cipher map $\fhat : \B^n \to \B^n$. A list
$[y_0, y_1, \ldots, y_{L-1}]$ of cipher values is represented by a
starting key $k_0 = k$ and the recurrence
\[
k_{i+1} = k_i \oplus h(y_i),
\qquad
y_i = \fhat(k_i),
\qquad i = 0, 1, 2, \ldots
\]
The iterator is the pair $(\,\text{current key } k_i,\ \fhat\,)$; its
two-line advance is \enquote{$y \leftarrow \fhat(k);\; k \leftarrow k
\oplus h(y)$.} Unrolling the recurrence, the $i$-th key is
\[
k_i \;=\; k \,\oplus\, \bigoplus_{j<i} h(y_j),
\]
an accumulator over the entire history of values produced so far. Two
consequences follow. First, the keys $k_0, k_1, k_2, \ldots$ are a
sequence of $n$-bit strings with no visible relation to one another:
each is the previous key XORed with the hash of an opaque value, so
under the random-oracle model on $h$ they are computationally
unlinkable, and the untrusted machine evaluating $\fhat$ sees it
applied to what looks like fresh, unrelated keys. Second, the chain is
genuinely sequential: $k_{i+1}$ cannot be formed without $y_i =
\fhat(k_i)$, so the untrusted machine cannot skip ahead, jump to the
middle, or reorder the traversal. It can only walk forward from a key
it has been handed.

\paragraph{Termination by sentinel.} The chain needs an end. Reserve a
distinguished latent value $\bot_{\mathrm{end}}$ and arrange that the
trusted machine's construction places it as $y_{L}$: when the iterator
decodes its current value to the sentinel, the list is exhausted. The
sentinel does double duty. It marks length without a separate length
field (the structure does not announce $L$ up front), and it absorbs
the failure mode of the hash chain: if a collision in $h$ or a cycle
in the induced key sequence would otherwise make the orbit branch or
loop, the trusted machine can detect this at construction time and
reposition the sentinel so that traversal halts cleanly before any
ambiguity. We make the random-oracle assumption that drives the
unlinkability, and the sentinel handling that tames its failure modes,
explicit in \Cref{sec:leakage}.

The essential point is the inversion of reading. \Cref{def:orbit}
calls $\{k_0, k_1, \ldots\}$ (equivalently $\{y_0, y_1, \ldots\}$) an
orbit and \Cref{thm:orbit-bound} treats its size as a leakage budget.
Here we have not measured an orbit; we have \emph{designed} one. The
shape of the orbit, a single forward chain of length $L$, is the shape
of the data structure, a list of length $L$. One cipher map and one
key encode the whole list.

\paragraph{Two-level orbits give maps.} A map is a list whose entries
are themselves iterable. Let the value space of $\fhat$ be (cipher
encodings of) key/list pairs: walking the outer orbit visits map keys
$\kappa_0, \kappa_1, \ldots$, and at each the decoded value carries the
starting key of an inner orbit encoding the list of values stored
under $\kappa_i$. The same $\fhat$ and the same advance rule serve both
levels; only the starting keys differ. The construction is uniform:
nest the orbit and you have nested the structure. A map is two levels,
a map of maps is three, and the recursion bottoms out at the list of
\Cref{def:orbit}.

\paragraph{Indexed access gives vectors.} A vector is a list whose
$i$-th element is reachable by a known offset. With the forward chain
above, reaching index $i$ costs $i$ advances: the structure is a
singly linked list, not a random-access array. One can shorten this by
storing, at selected stride points, the accumulated key $k_{i}$ as a
checkpoint (a skip list over the orbit), trading space for traversal
length; but the access remains sequential between checkpoints. Genuine
$O(1)$ random access is not available from a forward hash chain alone,
for the same structural reason the chain is unlinkable: forming $k_i$
requires the values before it. We return to this as the construction's
main open problem below.

\paragraph{Single cipher map, many data structures.} Across all of the
above, the cipher map $\fhat$ is fixed; what varies is the starting
key. Two lists $[y_0, y_1, \ldots]$ and $[y'_0, y'_1, \ldots]$ are the
orbits of $\fhat$ from two different keys $k$ and $k'$, and a map is a
family of such orbits indexed by an outer orbit, again all under the
same $\fhat$. The trusted machine therefore constructs \emph{one}
cipher map and obtains an unbounded supply of cipher data structures by
choosing where to start. This is the orbit analogue of the dispatch
pattern's economy: dispatch put many operations in one map by tagging
the input; designed orbits put many structures in one map by seeding
the iteration. In both, \enquote{data is code} is realized as a single
cipher map.

\begin{remark}[Open problem: random access]
\label{rem:random-access}
The natural completion of the vector construction is to make the index
a first-class cipher value. One would like a cipher map that takes a
cipher index $\enc(i, \cdot)$ together with the structure's starting
key directly to the $i$-th orbit element, with no $i$-step walk. This
is the dispatch pattern again, with the index in the role of the
operation tag: $\fhat(\enc(i,\cdot),\, k) = y_i$. The obstruction is
that the forward chain defines $y_i$ only through the values $y_0,
\ldots, y_{i-1}$, so a one-shot map from $i$ to $y_i$ must internalize
that dependence at construction time, which costs space growing with
the vector length and reintroduces, through the index's own orbit, the
leakage the chain was hiding. Whether a hash-chain-based cipher vector
admits oblivious random access within a useful space and leakage budget
is open. We state it as a problem, not a construction.
\end{remark}

%% ====================================================================
\section{Leakage of Orbit-Encoded Data Structures}
\label{sec:leakage}
%% ====================================================================

We now state precisely what an orbit-encoded list reveals to an
untrusted machine that holds $\fhat$ and $h$ and is given a starting
key, and is then free to walk the chain. We separate three quantities:
the length, the iteration order, and the per-element latent meanings.
The first two leak; the third is bounded by \Cref{thm:orbit-bound}. We
work under one assumption, stated once and used throughout.

\begin{assumption}[Random oracle]
\label{ass:ro}
The hash $h : \B^n \to \B^n$ is modeled as a random
oracle~\cite{bellare1993random}: its outputs are uniform and
independent across distinct inputs, and the untrusted machine learns
$h(y)$ only by querying $y$.
\end{assumption}

\paragraph{Length leaks, and how much.} As the untrusted machine walks
the chain it observes one new key per step until it decodes the
sentinel, so after a full traversal it knows $L$ exactly. The
\emph{a priori} leakage, before traversal, is the information the
length carries; we bound it in the natural way.

\begin{proposition}[Length leakage]
\label{prop:length-leak}
Let the list length $L$ be a random variable supported on $\{0, 1,
\ldots, n_{\max}\}$ with distribution induced by the trusted machine's
inputs. The length information an untrusted machine obtains by
traversing the orbit is
\[
I(L; \text{traversal}) \;=\; H(L) \;\leq\; \log_2 (n_{\max} + 1).
\]
In particular, for a list of length $n$ drawn from any distribution on
lengths up to $n$, the orbit leaks at most $\log_2 n + O(1)$ bits of
length information.
\end{proposition}

\begin{proof}
A complete traversal determines $L$, so the observation includes $L$
and $H(L \mid \text{traversal}) = 0$, whence $I(L; \text{traversal}) =
H(L)$. A random variable supported on a set of size $n_{\max}+1$ has
entropy at most $\log_2(n_{\max}+1)$, with equality for the uniform
distribution. The stated $\log_2 n + O(1)$ is this bound at $n_{\max} =
n$.
\end{proof}

This is a real cost and we do not hide it: an orbit-encoded list
announces its length to anyone who walks it. It is, however, a small
and structural cost, $O(\log n)$ bits regardless of what the $n$
elements mean, and it is the same length information a plaintext linked
list reveals through its node count. If even the length must be hidden,
the trusted machine pads with sentinel-suppressed decoy elements; we do
not develop padding here.

\paragraph{Order leaks.} The chain is sequential by construction
(\Cref{sec:designed-orbits}): $k_{i+1}$ requires $y_i$, so the
untrusted machine necessarily learns the iteration order, that element
$y_i$ precedes $y_{i+1}$. What it does \emph{not} learn from the key
sequence is any relation among the keys beyond \enquote{next}: under
\Cref{ass:ro} the keys $k_i = k \oplus \bigoplus_{j<i} h(y_j)$ are,
conditioned on the values being opaque, a sequence of strings each
offset from the previous by a uniform independent hash output, so the
key sequence alone is computationally unlinkable, revealing succession
but not, for instance, whether two different lists share a tail.
Iteration order is thus the exact order information leaked: a total
order on positions, and nothing finer.

\paragraph{Per-element meanings are bounded by the orbit.} The values
$y_0, \ldots, y_{L-1}$ are cipher values; their latent meanings are
what the data structure is for, and what its confidentiality must
protect. The traversal makes the orbit $\orbitF(k) = \{k_0, \ldots,
k_{L}\}$ (or equivalently the reached cipher values) available to the
untrusted machine, so each element's latent value $X_i$ is exactly in
the position of \Cref{thm:orbit-bound} with the operation set $F =
\{\text{advance}\}$ generating a chain of length $L+1$.

\begin{proposition}[Per-element confidentiality of an orbit list]
\label{prop:per-element}
Let $X_i$ be the latent value of the $i$-th element, drawn from the
trusted machine's per-element input distribution, and let the
untrusted machine's observation be the reached orbit of size $L+1$.
Then for each $i$,
\[
H(X_i \mid \mathcal{V}) \;\geq\; H(X_i) - \log_2 (L+1).
\]
\end{proposition}

\begin{proof}
The reached orbit has $L+1$ elements, so the untrusted machine's
observation of where element $i$'s computation resolves is supported on
a set of size $L+1$. \Cref{thm:orbit-bound} with $|\orbitF(k)| = L+1$
gives the bound directly.
\end{proof}

Two honest caveats accompany this bound, both inherited from
\Cref{thm:orbit-bound}. It is loose: it counts orbit elements and
ignores the functional ties among them (knowing $y_{i+1} =
\fhat(k_i \oplus h(y_i))$ constrains compatible latent values beyond
mere set size), so the true residual uncertainty is at least this
large and may be larger. And it degrades as the list grows: for $L$
comparable to $2^{H(X_i)}$ the per-element bound becomes vacuous, which
says that a long enough orbit, fully traversed, can in principle pin
down its elements. The construction's confidentiality is therefore
strongest for short-to-moderate lists and for elements drawn from
high-entropy domains, and the trusted machine controls it by the same
levers as any orbit: multiple representations per element fan the chain
out and shorten the effective per-step leakage, and capping $L$ caps
the bound.

\paragraph{What we do not claim.} It is tempting to assert that the
untrusted machine learns \emph{nothing} about the per-element meanings
beyond the orbit-size bound, i.e.\ that the chain hides the values up
to that bound. We do not prove this, and we flag it as a conjecture.

\begin{conjecture}[Per-element hiding]
\label{conj:hiding}
Under \Cref{ass:ro}, an untrusted machine that traverses an
orbit-encoded list of length $L$ learns the length and the iteration
order, and learns no more about the multiset of per-element latent
values $\{X_0, \ldots, X_{L-1}\}$ than is permitted by
\Cref{prop:per-element} applied per element, up to negligible
advantage.
\end{conjecture}

The conjecture is plausible because the keys are unlinkable and the
values are cipher values, but a proof would have to rule out
cross-element correlations introduced by the shared $\fhat$ and the
accumulator, which \Cref{thm:orbit-bound} explicitly does not control.
We leave it open rather than overclaim.

\paragraph{Collisions and cycles.} \Cref{ass:ro} makes a collision $h(y)
= h(y')$ for $y \neq y'$ occur with probability $\approx 2^{-n}$ per
pair, so over a list of length $L$ the chance any two advance steps
collide is $O(L^2 2^{-n})$, negligible for $n$ at cryptographic widths
and $L$ polynomial. When a collision does occur it can make the key
sequence revisit a previous key, turning the forward chain into a
\enquote{$\rho$}-shape with a cycle. The sentinel resolves this
without a runtime check on the untrusted side: the trusted machine,
which knows the values at construction time, detects any premature key
repetition and places the sentinel before it, so a well-formed
orbit-encoded list is collision-free by construction along its used
prefix. This converts a probabilistic hazard into a construction-time
obligation, and it is why the leakage statements above can speak of a
clean chain of length $L$.

\paragraph{The duality of orbits.} The same object that
\Cref{thm:orbit-bound} introduced as a leakage measure is here a
construction, and the bound it came with is exactly the tool that
quantifies the construction's leakage. The destructive reading (an
adversary explores an orbit and learns) and the constructive reading
(a designer shapes an orbit and stores) are two views of one
structure, and \Cref{prop:length-leak,prop:per-element} are the
constructive reading saying, in the designer's voice, what the
destructive reading always meant.

%% ====================================================================
\section{Code is Data: The Cipher Exponential}
\label{sec:exponential}
%% ====================================================================

The previous three sections realized data as code: a cipher data
structure is a cipher map, whether by dispatch or by designed orbit.
The reverse direction is already a theorem of the type algebra, and we
recall it to close the duality.

\begin{proposition}[Cipher exponential is the cipher map abstraction
{\cite[Prop.~4.3]{towell2026algebraic}}]
\label{prop:exponential}
For finite types $A$ and $B$, the cipher exponential $\cipher{A \to B}$
is precisely the set of cipher maps for functions $f : A \to B$. A
total function $\fhat : \B^n \to \B^n$ is a well-formed element of
$\cipher{A \to B}$ if and only if it satisfies the four cipher-map
properties of \Cref{def:cipher-map}.
\end{proposition}

The statement is small and its proof is a definitional unfolding (the
well-formedness conditions for membership in $\cipher{A \to B}$ are
exactly the four properties), so we do not reprove it. Its force here
is interpretive. A cipher map $\fhat$ wears two hats at once. As the
cipher functor's action on a morphism, $\fhat = \cipherS{f}{s}$ is
\emph{code}: the realization of the latent function $f$. As an element
of the cipher space $\cipherS{A \to B}{s}$, the very same bit-string
function $\fhat$ is a \emph{cipher value} of function type, a point of
a cipher type that some other cipher map could in turn take as input or
produce as output. Code is data, and the data in question is a cipher
value like any other, decodable (to the latent function) only by a
holder of the secret.

\paragraph{One carrier, both directions.} Put the two halves together.
\Cref{sec:dispatch,sec:designed-orbits} send a data structure to a
cipher map; \Cref{prop:exponential} identifies a cipher map with a
cipher value of function type. The cipher map is therefore the carrier
of both directions of the duality. A list is a cipher map (its iterator
is $\fhat$ seeded at a key); that cipher map is a cipher value of type
$\cipher{\B^n \to \B^n}$; and that cipher value could be stored as an
element of another orbit-encoded list, whose iterator is again a cipher
map. The regress does not bottom out in a separate notion of storage or
a separate notion of procedure, because in trapdoor computing there is
only one kind of thing flowing through the untrusted machine: bit
strings, carried by cipher maps. That is the precise sense in which the
cipher map is homoiconic. The categorical packaging, $\cipherS{X}{s}$ a
functor and $\cipherS{A \to B}{s}$ its exponential object, is the
bookkeeping that keeps the two readings of $\fhat$ consistent; it
labels the duality but does no work in establishing it, which rests
only on \Cref{prop:exponential} and the constructions of
\Cref{sec:dispatch,sec:designed-orbits}.

%% ====================================================================
\section{The Duality, Unified}
\label{sec:duality}
%% ====================================================================

Collecting the threads, the duality has three corners, and all three
are the same statement: in trapdoor computing, values, functions, data
structures, and the secret are cipher data carried by cipher maps.

\begin{center}
\begin{tabular}{lll}
\toprule
direction & construction & home \\
\midrule
data $\to$ code & data structure as a tagged cipher map / designed orbit & \Cref{sec:dispatch},~\ref{sec:designed-orbits} \\
code $\to$ data & cipher map as a point of $\cipher{A \to B}$ & \Cref{sec:exponential} \\
secret $\to$ data & cipher rekeying transforms the captured secret & \cite{towell2026rekeying} \\
\bottomrule
\end{tabular}
\end{center}

\paragraph{Two charted endpoints.} The code-to-data and
secret-to-data corners are already developed in the literature, and
they anchor the frame. The cipher exponential
(\Cref{prop:exponential}) is the code-to-data endpoint: it is settled,
a definitional identity. Cipher rekeying~\cite{towell2026rekeying} is
the secret-to-data endpoint, and it is the cleanest demonstration that
the secret is just captured state. A cipher object closed over a secret
$s_1$ is re-parameterized to one closed over $s_2$ by a rekeying map
$\rekey{s_1}{s_2} : \cipherS{X}{s_1} \to \cipherS{X}{s_2}$, which is
itself a cipher map the untrusted machine can apply, and which commutes
with every cipher operation (rekeying before, after, or between
operations gives the same result). Treating the secret as data is not a
metaphor: the operation that transforms it is an ordinary cipher map,
constructed by the trusted machine and run blindly by the untrusted
one, exactly like the operations of any cipher data structure. The
data-to-code corner, the dispatch pattern and designed orbits, is this
paper's contribution and completes the triangle: with all three corners
realized by cipher maps, the carrier is uniform across the entire
duality.

\paragraph{Mutability is captured cells.} A closure that captures a
mutable cell realizes a mutable data structure
(\Cref{ex:cipher-counter}); nothing in the framework forbids it. It is
worth being precise about what mutability costs in confidentiality, and
the answer is: nothing beyond the immutable version. A mutable
orbit-encoded structure that is advanced through states $c_0, c_1,
\ldots, c_t$ exposes, to an active untrusted machine, exactly the orbit
$\{c_0, c_1, \ldots, c_t\}$. The immutable \enquote{replay} of the same
history, the list $[c_0, c_1, \ldots, c_t]$ built once and traversed,
exposes the same set under the same advance map. So the orbit-closure
profile of a mutable closure equals that of its immutable replay, and
\Cref{thm:orbit-bound} and \Cref{prop:per-element} bound both
identically. In
trapdoor computing, mutation does not leak more than recording the
sequence of states would; it is the orbit that leaks, not the
mutability. This is what licenses treating mutable and immutable cipher
data structures within one leakage theory.

\paragraph{The thesis, in one line.} There is one carrier. A cipher
value is a point of a cipher space; a cipher map is both code and a
point of a cipher exponential; a cipher data structure is a cipher map
by dispatch or designed orbit; and the secret is captured state that
rekeying moves. All four are cipher data, and the cipher map carries
them all.

%% ====================================================================
\section{Related Work}
\label{sec:related}
%% ====================================================================

\paragraph{Homoiconicity and closures as data.} That code and data can
be the same object is as old as Lisp~\cite{mccarthy1960recursive}, and
the closure as the device that fuses them is the lesson of Landin's
SECD machine~\cite{landin1964mechanical} and of Abelson and Sussman's
data-as-procedures
exposition~\cite{abelson1996sicp}, whose $\mathrm{cons}/\mathrm{car}/
\mathrm{cdr}$ we used throughout. Our contribution is not to rediscover
homoiconicity but to instantiate it over a secret: the cipher closure
is the homoiconic object of trapdoor computing, and the cipher map is
the single carrier in which the duality is realized. Where SICP's pair
is an arbitrary procedure, ours must be a cipher map (a total function
with the four properties), and that constraint is what makes the
information-theoretic leakage analysis both necessary and possible.

\paragraph{Oblivious data structures and ORAM.} Oblivious
RAM~\cite{goldreich1996software} and modern constructions such as Path
ORAM~\cite{stefanov2013path}, and the oblivious data structures built
on them~\cite{wang2014oblivious}, address a different problem on a
different axis. They assume the data lives in untrusted storage and
hide the \emph{access pattern}: which location is touched, in what
order, so that two access sequences of the same length are
indistinguishable. The data structure is conventional; the protocol
that reaches into it is what is made oblivious, typically at
polylogarithmic bandwidth overhead per access. Designed orbits make a
different move. The data structure is not stored and accessed; it
\emph{is} a cipher map, and there is no access pattern over storage to
hide because there is no storage, only the evaluation of a total
function on bit strings. What we analyze instead is what the cipher
map's operation interface leaks information-theoretically, the length
and iteration order of \Cref{sec:leakage}, bounded by orbit closure.
The two approaches are orthogonal and composable: one could run an
orbit-encoded cipher data structure's iterator over ORAM-protected
storage to hide both the access pattern and the latent values, but
neither subsumes the other, and our guarantees are information-theoretic
where ORAM's are typically computational.

\paragraph{Functional data structures.} The persistent, purely
functional data structures catalogued by Okasaki~\cite{okasaki1998purely}
share our reliance on immutable structure and on representing data as
the closures that build it, and the orbit-encoded list is, in shape, a
functional cons-list whose tail is reached by advancing the hash chain.
The concerns differ: Okasaki optimizes amortized and worst-case time
under persistence, whereas we are concerned with what the structure
reveals to an untrusted evaluator. The connection is suggestive (a
persistent structure's sharing of tails is exactly the cross-list tail
correlation that \Cref{sec:leakage} notes the unlinkable key sequence
hides), and a confidentiality-aware account of standard functional data
structures over cipher maps is natural future work.

\paragraph{Quantitative information flow.} Our leakage statements are
in the QIF idiom: confidentiality measured as residual uncertainty
about a secret given an observation, with mutual information and
conditional entropy as the currency~\cite{shannon1948mathematical}, and
with the now-standard refinements (min-entropy and $g$-leakage) that
sharpen Shannon leakage into operational
guarantees~\cite{smith2009foundations}. The orbit-closure
bound~\cite{towell2026algebraic} we use is a counting bound in exactly
this tradition. We have not pursued the min-entropy or $g$-leakage
forms of the orbit bound for designed orbits; doing so, and tightening
the loose counting bound of \Cref{thm:orbit-bound} for the specific
advance operation of a hash chain, is open (and connects to
\Cref{conj:hiding}).

%% ====================================================================
\section{Conclusion}
\label{sec:conclusion}
%% ====================================================================

In trapdoor computing there is one kind of object on the untrusted
machine: a total function on bit strings, the cipher map, applied to
opaque cipher values. We have argued that this single carrier is enough
to realize the code--data duality of the Lisp tradition, with the
cipher closure as its homoiconic object. Three results make the
duality precise. A cipher data structure is a cipher map, either by the
dispatch pattern, which folds an operation interface into a query tag,
or by a designed orbit, which folds sequential structure into the shape
of a hash chain so that one cipher map and a starting key encode a
list, a map, or a vector (\enquote{data is code}). A cipher map is, by
the cipher exponential identity, a cipher value of function type
(\enquote{code is data}). And the captured secret is itself state that
cipher rekeying transforms by an ordinary cipher map (\enquote{secret
is data}). The same orbit that the type algebra introduced as a leakage
budget is here a construction primitive, and the bound that limited the
adversary is the bound that measures, honestly, what an orbit-encoded
data structure reveals: its length, its iteration order, and per-element
meanings only up to the orbit-size bound.

We have been careful to separate what is proved from what is not. The
length and order leakage are established (\Cref{prop:length-leak} and
the sequentiality of the chain); the per-element bound follows directly
from orbit closure (\Cref{prop:per-element}); but the stronger claim
that the construction hides everything else up to that bound is a
conjecture (\Cref{conj:hiding}), because the shared cipher map and the
accumulator introduce cross-element correlations the counting bound
does not control. The orbit bound itself is loose, and tightening it
for the hash-chain advance is open.

Three directions stand out for future work. The first is oblivious
random access for the cipher vector (\Cref{rem:random-access}): making
the index a cipher value and dispatching on it without an $i$-step walk,
within a useful space and leakage budget, or proving it cannot be done.
The second is mutable orbit reshaping: \Cref{sec:duality} showed that a
mutable closure's leakage equals its immutable replay's, but structural
edits that re-shape an orbit in place (insertion, deletion, splicing)
deserve their own leakage account. The third, and the most
far-reaching, is the limit of \enquote{data is code}: a cipher Turing
machine, in which the structure encoded by an orbit is a computation
rather than a list, would make code a fully general cipher data
structure and close the loop with the program-realization line of
work~\cite{towell2026cipherprog}. In each case the discipline is the
same one this paper has tried to keep: build the construction as
mathematics over a one-way hash and a uniform representation, measure
its leakage information-theoretically, and claim no more than the orbit
allows.

\bibliographystyle{plainnat}
\bibliography{references}

\end{document}
